\section{Unconditional source identities and complete row estimates}
\label{sec:source-dependencies}

This section proves the entry estimates used below. All prime sums contain
all prime powers. No conjectural assertion about the critical-line location or simplicity
of the nontrivial zeros, or about the sign of an explicit-formula quadratic
form, is used.

\subsection{Normalization, the pole, and the ground row}
Let
\[
 \xi(s)=\tfrac12s(s-1)\pi^{-s/2}\Gamma(s/2)\zeta(s),\qquad
 \Xi(z)=\xi(\tfrac12+iz),\qquad
 \widehat u(z)=\int_{\R}u(x)e^{-izx}\,\dd x.
\]
Jacobi inversion gives the smooth even continuation of the displayed
positive theta series for $\kappa$. For completeness, if
$\vartheta(t)=\sum_{j\geq1}e^{-\pi j^2t}$, that series, with $t=e^{2x}$,
is
\[
 \kappa(x)=e^{x/2}\{4t^2\vartheta''(t)+6t\vartheta'(t)\}.
\]
Jacobi inversion cancels the leading $t^{-1/2}$ term in the braces at zero.
Mellin substitution and two integrations by parts, first for $\Re s>1$,
therefore give
\[
 \int_{\R}\kappa(x)e^{(s-1/2)x}\,\dd x
 =\frac{s(s-1)}2\int_0^\infty t^{s/2-1}\vartheta(t)\,\dd t
 =\xi(s).
\]
Both sides continue to entire functions. In particular,
\begin{equation}\label{eq:dep-normalization}
 \widehat\kappa=\Xi,\qquad
 \int_{\R}2C(x)\kappa(x)\,\dd x=1.
\end{equation}
These are classical theta--Mellin identities; see \cite{Romik2021}.

Write
\[
 D_\Gamma v(x)=\int_0^\infty\rho(s)
       \{2v(x)-v(x+s)-v(x-s)\}\,\dd s,
 \qquad
 c_\Gamma=\log\pi+\gamma+\frac\pi2+3\log2.
\]
The difference inside the integral is kept coupled at zero. The complete
source row is
\begin{equation}\label{eq:dep-ground-row}
 P\kappa=C+(D_\Gamma-c_\Gamma)\kappa,
 \qquad P=P_-+P_+.
\end{equation}
We record explicitly its unconditional origin, including the pole.
For real even compact smooth physical tests put
\[
 \mathcal Q(u)=\int_0^\infty\rho(s)
             \|u(\cdot+s)-u\|_{L^2(\R)}^2\,\dd s
 -c_\Gamma\|u\|_{L^2(\R)}^2
 -2\sum_{n\ge2}\frac{\Lambda(n)}{\sqrt n}
                      \langle u,u(\cdot+\log n)\rangle
 +2\left|\int_{\R}Cu\right|^2.
\]
The polarized classical explicit formula, in this normalization, is
\begin{equation}\label{eq:dep-explicit}
 \mathcal Q(u,v)=\sum_{\rho_\zeta}
    \overline{\widehat u(\overline{z_{\rho_\zeta}})}
                        \widehat v(z_{\rho_\zeta}),
 \qquad z_{\rho_\zeta}=(\rho_\zeta-\tfrac12)/i,
\end{equation}
with all nontrivial zeros and their multiplicities. The archimedean
normalization and the two pole terms are those of
\cite[Appendix B]{ConnesConsani2021}. More explicitly, the archimedean
expression before rewriting is
\[
 -\{\log(4\pi)+\gamma\}\|u\|_2^2
 -2\int_0^\infty
 \frac{e^{s/2}\langle u,u(\cdot+s)\rangle-\|u\|_2^2}
      {e^s-e^{-s}}\,\dd s.
\]
The identity
$2\int_0^\infty(e^{s/2}-1)/(e^s-e^{-s})\,\dd s
=\pi/2+\log2$ gives exactly the displayed $D_\Gamma,c_\Gamma$ convention.

Here is the passage to $u=\kappa$, rather than a formal substitution.
Every fixed derivative of $\kappa$ decreases faster than every exponential.
For a smooth physical function define
\[
 S_2(u)=\sum_{j=0}^2\int_{\R}e^{|x|}|u^{(j)}(x)|\,\dd x.
\]
Twice integrating by parts bounds $|\widehat u(z)|$ by a constant times
$S_2(u)/(1+|z|)^2$ on $|\Im z|\le1/2$. The unconditional zero-counting
bound $N(T)=O(T\log(2+T))$ makes the polarized sum continuous under
$S_2$ convergence of smooth cutoffs. The physical terms are continuous
under $H^1$ convergence together with convergence of
$e^{|x|}u$ in $L^2$: indeed
\[
 |\langle u,v(\cdot+s)\rangle|
 \le e^{-|s|}\|e^{|x|}u\|_2\|e^{|x|}v\|_2,
 \qquad \sum_{n\ge2}\Lambda(n)n^{-3/2}<\infty.
\]
The Gamma term is controlled by its small-shift derivative bound and its
integrable large-shift tail. Theta cutoffs have all these convergences.
Since $\Xi$ vanishes at every $z_{\rho_\zeta}$ and its conjugate,
\eqref{eq:dep-explicit} gives $\mathcal Q(\kappa,v)=0$ for every compact
smooth even $v$. The pole term is
$2(\int C\kappa)(\int Cv)=\int Cv$, by
\eqref{eq:dep-normalization}. Thus the distributional identity is
$D_\Gamma\kappa-c_\Gamma\kappa-P\kappa+C=0$.
The local continuity proved next makes it the pointwise identity
\eqref{eq:dep-ground-row}. In particular, this argument does not infer
radicality from a zero quadratic value in an indefinite form.

\subsection{Derivative tails and a positive derivative source}
On $x\ge0$ put $t_j=\pi j^2e^{2x}$ and $h_0(t)=4t^2-6t$.
Differentiation of a theta summand acts on its polynomial by
\[
 \mathscr Lh=(\tfrac12-2t)h+2th',\qquad h_k=\mathscr L^kh_0.
\]
For any fixed $k$, $h_k$ has degree $k+2$. A coefficient-sum bound,
$t^q e^{-t}\le(2q)^q e^{-t/2}$, and
$j^2\ge1+3(j-1)$ give constants $A_k$ such that
\begin{equation}\label{eq:dep-derivative-tail}
 |\kappa^{(k)}(x)|\le
 A_ke^{|x|/2}\frac{\exp[-(\pi/2)e^{2|x|}]}
                         {1-e^{-3\pi/2}}.
\end{equation}
Parity supplies negative $x$. This bounds the full series, not just its
first term. Translated derivatives at $y\pm\log n$ consequently have
Gaussian-in-$n$ tails, uniformly on compact $y$ intervals, even after
multiplication by $\Lambda(n)/\sqrt n\le\log n/\sqrt n$.
The differentiated Gamma integrands are $O(s^2)$ near zero and are
integrable at infinity. Thus the complete rows may be differentiated
twice, and integration by parts against $C$ has no boundary terms.

Define
\[
 r=\frac{\kappa''+20\kappa}{47},\quad
 g_0=\kappa-r=\frac{27\kappa-\kappa''}{47},\quad m=\frac{81}{188}.
\]
\begin{lemma}\label{lem:dep-positive-source}
The following inequalities hold pointwise on the whole line:
\[
 r\ge\kappa/188>0,\qquad r\ge\kappa\mathbf1_{\R\setminus I},
 \qquad g_0\le\kappa\mathbf1_I,\qquad r''+40r\ge0.
\]
Moreover,
\begin{align}
 \int_{\R}2Cr\,\dd x&=m,\label{eq:dep-r-mass}\\
 Pr&=mC+(D_\Gamma-c_\Gamma)r,\label{eq:dep-r-row}\\
 Pg_0&=\frac{107}{188}C+(D_\Gamma-c_\Gamma)g_0.
 \label{eq:dep-g-row}
\end{align}
\end{lemma}
\begin{proof}
Exact differentiation gives
\[
 h_2=16t^4-112t^3+165t^2-\tfrac{75}{2}t,
\]
\[
 h_4=64t^6-1056t^5+5176t^4-8512t^3
                 +\tfrac{15465}{4}t^2-\tfrac{1875}{8}t.
\]
The ratio
\[
 R_2(t)=h_2(t)/h_0(t)=4t^2-22t+\tfrac{33}{4}+\frac6{2t-3}
\]
satisfies $R_2'(t)\ge2/3$ for $t\ge3$,
$R_2(3)=-79/4$, and $R_2(25/4)=525/19$.
All theta weights are positive, with $t_j>3$. Hence
$r/\kappa\ge(1/4)/47=1/188$.
On $|x|\ge a$, $t_j\ge2\pi>25/4$, so
$r/\kappa\ge(525/19+20)/47=905/893>1$.

For the last differential inequality,
\[
 h_4+60h_2+800h_0=(t/8)Q(t),
\]
\[
 Q(t)=512t^5-8448t^4+49088t^3-121856t^2+135730t-58275.
\]
On the three intervals $[j,j+1]$, $j=3,4,5$, write
$5Q(j+u)=\sum_{i=0}^5 B_{j,i}\binom5i u^i(1-u)^{5-i}$.
The exact coefficient rows are
\[
\begin{array}{c|rrrrrr}
 j& B_{j,0}&B_{j,1}&B_{j,2}&B_{j,3}&B_{j,4}&B_{j,5}\\\hline
3&88575&113521&139459&163285&181127&190905\\
4&190905&200683&202397&193967&175105&149875\\
5&149875&124645&93047&59145&31355&25005
\end{array}
\]
and are positive. For $u\ge0$,
\[
 Q(6+u)=512u^5+6912u^4+30656u^3+42880u^2-6350u+5001.
\]
The final quadratic is positive since
$4\cdot42880\cdot5001-6350^2=817449020>0$.
Thus $Q>0$ on $[3,\infty)$, and summation proves $r''+40r\ge0$.

Two integrations by parts give $\int2C\kappa''=1/4$, so
\eqref{eq:dep-r-mass} follows from \eqref{eq:dep-normalization}.
Applying $(\partial_x^2+20)/47$ to \eqref{eq:dep-ground-row}, using
$C''=C/4$ and the convergence just proved, gives
\eqref{eq:dep-r-row}. Subtraction gives \eqref{eq:dep-g-row}.
\end{proof}

\subsection{A complete Gamma correction and weighted rows}
\begin{lemma}\label{lem:dep-gamma-correction}
If $v\ge0$ is $C^2$, $v''+40v\ge0$, and its compensated Gamma row exists,
then
\[
 (D_\Gamma-c_\Gamma)v\le-\frac{37}{350}v.
\]
Consequently,
\begin{equation}\label{eq:dep-upper-rows}
 Pr\le mC-\frac{37}{350}r\le mC,\qquad
 P\kappa\le C-\frac{37}{350}\kappa<C,
 \qquad q_P:=\frac{P\kappa}{2C}<\frac12.
\end{equation}
Also
\begin{equation}\label{eq:dep-weighted-source}
 \kappa(x)e^{|x|}\le282r(x),\qquad
 B_L:=P_-(\kappa e^{|\cdot|})\le282mC.
\end{equation}
\end{lemma}
\begin{proof}
Put $k=\sqrt{40}$ and
$F_x(s)=v(x+s)+v(x-s)-2v(x)\cos(ks)$.
Then $F_x(0)=F_x'(0)=0$ and $F_x''+40F_x\ge0$.
Variation of constants, whose sine kernel is nonnegative for
$0\le s\le\pi/k$, gives $F_x\ge0$ there. In particular, on
$0\le s\le1/4$,
\[
 2v(x)-v(x+s)-v(x-s)\le40s^2v(x).
\]
For larger $s$ use $2v(x)$. Since
$\rho(s)=e^{s/2}/(2\sinh s)<4/(7s)$ on $(0,1/4]$,
\[
 40\int_0^{1/4}s^2\rho(s)\,\dd s<5/7.
\]
With $u=e^{-1/8}<8/9$ the exact antiderivative gives
\[
 \int_{1/4}^\infty\rho(s)\,\dd s=\operatorname{atanh}u+\arctan u,
\]
so $D_\Gamma v\le(5/7+\log17+\pi/2)v$.
Now $\gamma>H_6-\log7>1/2$, because $H_6=49/20$ and
$\log7<39/20$. Also $8\pi>25$ and
$\log(25/17)>8/25$ (integrate $1/t$ over $[17,25]$).
Therefore
\[
 c_\Gamma-(5/7+\log17+\pi/2)
 >\tfrac12+\tfrac8{25}-\tfrac57=\frac{37}{350}.
\]
This proves the comparison. It applies to $r$ by
Lemma~\ref{lem:dep-positive-source}, and to $\kappa$ because
$\kappa''+20\kappa\ge0$ and $\kappa>0$.
Equations~\eqref{eq:dep-ground-row} and \eqref{eq:dep-r-row}
prove \eqref{eq:dep-upper-rows}.

For $|x|\le a$, $r\ge\kappa/188$ and $e^{|x|}\le\sqrt2<3/2$.
For $x\ge a$, the same positive theta weights give
\[
 \frac r\kappa\ge\frac{A(\pi e^{2x})}{47},\qquad
 A(t)=R_2(t)+20=4t^2-22t+\tfrac{113}{4}+\frac6{2t-3}.
\]
For $t\ge25/4$, $A(t)>t^2$: the polynomial part of
$A(t)-t^2$ is increasing there and equals $127/16$ at $25/4$.
Hence $r/\kappa>(9/47)e^{4x}$.
Evenness supplies negative $x$, and the global comparison
$\kappa e^{|x|}\le282r$ follows. Positivity of the complete row gives
$B_L\le282Pr\le282mC$.
\end{proof}

\subsection{The core mass}
\begin{lemma}\label{lem:dep-core-mass}
For $a=\log2/2$ and $I=[-a,a]$,
\[
 K_I:=\int_I\kappa(x)\,\dd x>\frac{5780}{14391}>\frac25.
\]
\end{lemma}
\begin{proof}
For $h_1=\mathscr Lh_0=-8t^3+30t^2-15t$,
\[
 h_1+7h_0=-t(8t^2-58t+57).
\]
The quadratic is positive and increasing for $t\ge25/4$; its value
there is $7$. Thus $\kappa'+7\kappa<0$ for $x\ge a$ and
$\kappa(x)\le\kappa(a)e^{-7(x-a)}$.
At $a$, the first theta term is less than
\[
 \frac65\frac{475}{4}e^{-25/4}<\frac{684}{2465},
\]
since $e^{25/4}>400(1+1/4+1/32+1/384)=12325/24$.
The complete remaining tail is bounded by
\[
 \frac{1536}{5}\sum_{j\ge2}j^4 20^{-2j^2}
 <\frac{49152}{5\cdot20^8}<10^{-6}.
\]
Its successive ratio is below $1/2$. Since
$5/18-684/2465=13/44370>10^{-6}$, $\kappa(a)<5/18$.
With $e^{a/2}<6/5$ we obtain
\[
 \int_{\R\setminus I}2C\kappa
 <\frac59\left(\frac{6/5}{13/2}+\frac1{15/2}\right)
 =\frac{62}{351}.
\]
Moreover
$C(a)^2=1/2+3/(4\sqrt2)<577/560<(41/40)^2$, using
$\sqrt2>140/99$. Equation~\eqref{eq:dep-normalization} yields
\[
 \frac{41}{20}K_I>\int_I2C\kappa>\frac{289}{351},
\]
which is the claimed bound.
\end{proof}

\subsection{A uniform returning-core lower row}
\begin{lemma}\label{lem:dep-lower-left}
Set $q_I=(2C)^{-1}P(\kappa\mathbf1_I)$. Then
\[
 q_I(y)>1/16\quad (y\ge a),\qquad
 L(y)=P_-\kappa(y)\ge C(y)/8\quad (y\ge a).
\]
In particular this lower bound for $L$ is asserted only on the indicated
half-line, not on all positive hubs.
\end{lemma}
\begin{proof}
First, the positive first theta term
$F(x)=e^{x/2}h_0(\pi e^{2x})e^{-\pi e^{2x}}$, $x\ge0$,
has no interior minimum on $[0,b]$ for any $b$: the derivative polynomial
$h_1(t)=t(-8t^2+30t-15)$ has at most one sign change on $t\ge\pi$,
from positive to negative. Endpoint bounds therefore give
\begin{equation}\label{eq:dep-local-lower}
 \kappa(x)>21/80\ (|x|\le a),\quad
 \kappa(x)>1/3\ (|x|\le3/10),\quad
 \kappa(x)>2/5\ (|x|\le13/50).
\end{equation}
Here are explicit endpoint certificates. At zero,
$F(0)>(325/16)/(63/2)=325/504>2/5$.
At $a$, $h_0(2\pi)>475/4$, $2^{1/4}>19/16$ and
$e^{2\pi}<e^{44/7}<537$, giving $F(a)>21/80$.
At $13/50$, $e^{13/25}<17/10$ implies
$\pi e^{13/25}<27/5$. Since $h_0(t)e^{-t}$ decreases for $t\ge3$,
\[
 F(13/50)>\frac{113}{100}\frac{2106/25}{225}>2/5.
\]
We used $e^{13/100}>113/100$ and $e^{27/5}<225$.
Similarly, $e^{3/5}<11/6$, $\pi e^{3/5}<29/5$,
$e^{3/20}>23/20$ and $e^{29/5}<333$ yield
$F(3/10)>(23/20)(2494/25)/333>1/3$.

Let $a_2=\log2/\sqrt2>49/100$. On $[a,2/5]$, the $n=2$ branch
returns to $I$ with displacement at most $a$, and
$d(y)=1/(2C(y))>49/100$. Thus
\[
 q_I>(49/100)^2(21/80)>1/16.
\]
On $[2/5,9/20]$ its displacement is below $3/10$, while $d>2/5$,
so $q_I>(2/5)(49/100)(1/3)=49/750>1/16$.
The four closed intervals
\[
 [\log p-13/50,\log p+13/50],\qquad p=2,3,5,7,
\]
cover $[9/20,2]$. Indeed $\log2<69315/100000$ places the first
left endpoint below $9/20$; the largest adjacent prime ratio is $5/3$
and $e^{13/25}>5/3$; and $\log7>7/4$ places the last right endpoint
above $2$. Since $13/50<a$, each chosen branch genuinely returns to $I$.
On this interval $d>8/25$ because $\cosh1<25/16$.
Also $\log p/\sqrt p\ge a_2$: the function $\log t/\sqrt t$ increases
on $[2,7]$, as $e^2>7$. Thus
\[
 q_I>(8/25)(49/100)(2/5)=196/3125>1/16.
\]
Additional primes and prime powers only increase these lower bounds.

It remains to cover every $y\ge2$. The signed minorant $g_0$ from
Lemma~\ref{lem:dep-positive-source} obeys
\begin{equation}\label{eq:dep-g-tail}
 |g_0(x)|,\ |g_0''(x)|<\varepsilon:=1/100\qquad(|x|\ge3/2).
\end{equation}
For an explicit full-tail proof, at $t\ge60$ coefficient-sum bounds give
$|h_2|\le18t^4$, $|h_4|\le84t^6$, and
\[
 |27h_0-h_2|/47\le2t^6,\qquad
 |27h_2-h_4|/47\le2t^6.
\]
For $x\ge3/2$, $t_j>60j^2$ and $e^{x/2}\le t_j$.
The bound $t^7e^{-t}\le14^7e^{-t/2}$ and the full geometric sum give
both quantities in \eqref{eq:dep-g-tail} an upper bound
\[
 \frac{2\cdot14^7e^{-30}}{1-e^{-90}}
 <\frac{2\cdot14^7}{20^{10}(1-20^{-30})}<1/100.
\]
Evenness treats $x<0$.

For $y\ge2$, keep the Gamma difference coupled on $0<s\le1/2$.
Taylor's integral remainder gives
$2g_0(y)-g_0(y+s)-g_0(y-s)\ge-\varepsilon s^2$.
Since $\rho(s)\le2/(3s)$ there, its integrated cost is
$\varepsilon/12$. For $s\ge1/2$, the inequalities
$g_0\le\kappa\mathbf1_I$ and $g_0(y)\ge-\varepsilon$ give
\[
 2g_0(y)-g_0(y+s)-g_0(y-s)
 \ge-2\varepsilon-\kappa(y-s)\mathbf1_I(y-s).
\]
Here $\int_{1/2}^\infty\rho<4$; the returning-core term is at most
$1/2$, because its shifts exceed $3/2$, where $\rho<1$, and
$\int_I\kappa\le\tfrac12\int2C\kappa=1/2$.
The negative $2g_0(y)$ contribution has thus been retained, and
\[
 D_\Gamma g_0(y)\ge-\varepsilon/12-8\varepsilon-1/2.
\]
As $g_0(y)\le0$, $c_\Gamma>0$, and $d(y)\le1/3$ on $y\ge2$,
positivity of the complete prime row and \eqref{eq:dep-g-row} yield
\[
 q_I\ge dPg_0
 \ge\frac{107}{376}-\frac16-\frac{97}{3600}
 =\frac{15391}{169200}>\frac1{16}.
\]
Finally, for $y\ge a$ no $y+\log n$ lies in $I$. All returning-core
branches therefore point left, and
$L\ge P(\kappa\mathbf1_I)=2Cq_I\ge C/8$.
\end{proof}

\subsection{An outward weighted row and scalar certificates}
The other weighted row needed in the hub estimate also has a complete
analytic bound. By evenness, for $t\ge0$,
\[
 \frac{B_L(-t)}{L(-t)}
 =e^t\frac{\sum_{n\ge2}n a_n\kappa(t+\log n)}{R(t)}<3e^t.
\]
To check the last inequality, subtract the exact $n=2$ contribution
$2R$ before bounding the remaining numerator. Put $T=\pi e^{2t}>3$
and $\phi(u)=(4u^2-6u)e^{-u}$. The positive $n=2,j=1$ denominator,
$\phi(4T)\ge56T^2e^{-4T}$, and
$\phi(Tj^2n^2)\le4T^2j^4n^4e^{-Tj^2n^2}$ give
\[
 \frac{\sum n a_n\kappa(t+\log n)}{R(t)}
 \le2+\frac17\sum_{n\ge3,j\ge1}j^4n^6
                         20^{-(j^2n^2-4)}<3.
\]
Indeed the $j$ ratio is at most $16/20^{27}<1/2$ and the remaining
$n$ ratio is at most $(4/3)^6/20^7<1/2$; the double sum is at most
$4\cdot3^6/20^5=729/800000<1$.
With $h_L=W(R/L)B_L$ this implies, using $dR<1/2$,
\begin{equation}\label{eq:dep-h-rows}
 h_L(-t)<\tfrac32e^t\kappa(t)\quad(t\ge0),\qquad
 h_L(y)\le1128m\kappa(y)=486\kappa(y)\quad(y\ge a).
\end{equation}
For the second inequality combine $B_L\le282mC$, $L\ge C/8$,
$R<C$, and $W=\kappa/(2C)$. Reflection gives $h_R(y)=h_L(-y)$.

All auxiliary scalar brackets above can be checked by finite rational
series, without evaluating a theta or prime sum. A convenient general
certificate, for rational $x>0$ and $x<n+2$, is
\[
 E_n(x):=\sum_{j=0}^n\frac{x^j}{j!}<e^x
 \le U_n(x):=E_n(x)+\frac{x^{n+1}}{(n+1)!}
                         \frac1{1-x/(n+2)}.
\]
For the less familiar upper brackets used in this section, direct
rational substitution gives
\[
\begin{array}{c|c|c}
 x&n&U_n(x)\text{ is less than}\\\hline
44/7&50&537\\
13/25&12&17/10\\
27/5&40&225\\
3/5&12&11/6\\
29/5&40&333\\
1/3&12&3/2\\
1/4&12&4/3\\
1/8&12&8/7\\
2&25&8
\end{array}
\]
The lower brackets $E_8(13/25)>5/3$, $E_6(39/20)>7$,
$E_4(2)=7$, $E_3(3/4)>2$, and $E_4(3/2)>4$ are likewise rational.
In particular $e^{7/4}=e^2/e^{1/4}<8/(5/4)=32/5<7$.
The positive series
\[
 \log2=2\sum_{j\ge0}\frac1{(2j+1)3^{2j+1}}
\]
with its geometric remainder gives
$693/1000<\log2<69315/100000$ after seven terms.
Together with $\sqrt2<99/70$ it proves $a_2>49/100$.
Finally, positive even Taylor sums and their geometric remainders give
$\cosh(1/5)<50/49$, $\cosh(1/2)<8/7$, and $\cosh1<25/16$.
For example the first and last follow from
\[
 \sum_{j=0}^4\frac{(1/5)^{2j}}{(2j)!}
 +\frac{(1/5)^{10}}{10!}\frac1{1-1/(25\cdot11\cdot12)}<50/49
\]
and
\[
 \sum_{j=0}^6\frac1{(2j)!}
 +\frac1{14!}\frac1{1-1/(15\cdot16)}<25/16.
\]
These certificates, the elementary $25/8<\pi<22/7$, and rational
squaring establish all constants in this section independently of any
source samples or numerical spectral premises.
